
% ============================================================================
\clearpage
\section*{v4.10.0-r55 Source-Only: How the Open Problems Hang Together (Root, Cross-Sections, Remainders)}
\addcontentsline{toc}{section}{v4.10.0-r55 Source-Only: How the Open Problems Hang Together}
\label{sec:v4100-r55-cross-map}
% ============================================================================

\paragraph{Purpose.}
At the author's request (``how do all open questions in the paper hang together'', then ``record it'') an exploration of 23 November 2026
mapped the open problems onto one picture; at the author's later request it becomes a layer, so that the German article can refer to it.  The
layer updates the map with layers r53 and r54.  It contains no computation and no hypothesis; every statement is read from the cited layers and
notes.  Earlier layers are not rewritten.

\noindent\textbf{Source-only status.}
This v4.10.0-r55 layer is source-only.  It does not build a PDF, does not assign a new DOI,
and does not authorize a public release.  The v5.3.0 Version DOI
\href{https://doi.org/10.5281/zenodo.24434403}{10.5281/zenodo.24434403} and the Concept DOI
\href{https://doi.org/10.5281/zenodo.19210728}{10.5281/zenodo.19210728} are unchanged by it.

\subsection*{v4.10.0-r55 results}
\label{subsec:v4100-r55-results}

\paragraph{Remark v4.10.0-r55.1 (the root and its translations).}
Every value that OP2--OP5 supply today is a function of one number, \(\beta_{\rm sym}\) (equivalently \(\eta\), \(\gamma_L\), \(\gamma_M\)), which enters through five separately
set identifications: the spectral offset of OP2, the residual bend \(\pi/4-\gamma_M/2\) of OP3, \(|U_{e2}|^2=\beta^2\) and \(\varepsilon=\eta\) of OP4, \(\Omega_M=\beta^2\) of OP5, and
\(\cot\theta_W=\gamma_L\) (r30.6).  Since the unitary boost, the Holst duality rotation (r31.3, r31.4) and the frame map \(J_B\) (r52) supply no frame-dependent area
quantum, OP1 rests on the counting rules alone, U2 for M and \(j=\tfrac12\) for T.  A solution of OP1 would supply the number but none of the five
identifications.

\paragraph{Remark v4.10.0-r55.2 (cross-section Q1: the readout).}
Three places need the same kind of datum, a frame \(u\), an axis \(n\) and a time orientation: OP1 in reading B (\(\varphi=\operatorname{sech}2\eta\) relative to the field rest
frame \(R\)), OP3 (the sign of the octant step is a time orientation of the readout, r31.1) and OP5 (the directed coupling, r40, r45.3, r46).  r45.4 named
this for OP5 together with OP1 and OP3.  The cross-section splits into two questions.  Q1a, the orientation, concerns the octant sign (r31.1) and
the counting of OP1, whose weight two needs a readout that distinguishes the two orientations of a puncture without selecting one (r53.2, r53.4);
it is open.  Q1b, frame and axis, concerns OP1 in reading B and the readout branch of OP5; both readouts are quadratic in \(u\) and blind to the time
orientation (r48.2).  r54 answers Q1b negatively: in OP1 the field selects frame and axis as a minimum of the readout, in OP5 the rest frame of the
preparation is a saddle and frame and axis are an external setting of the detector.  A common law for frame and axis needs an additional setting.

\paragraph{Remark v4.10.0-r55.3 (cross-section Q2: the dynamics).}
The common action W1/W2 is frozen (r05); it fixes no scale for \(\Lambda\), for the rate of pair creation or for the duration of the transition at the big bang.
No candidate meets the mediator profile A1--A5 (r30.7), and the directed mediator of OP5 is absent from the record (r34--r49; vacuum excluded by
r49--r51).  The deferred idea I1, a string between M and T, is the only candidate that touches OP1 (T counts string endpoints), OP3 (one crossing,
CT-002) and OP5 (r37 comes closest to the mediator term) at once.  Loosely, the open question on unequal string relaxation of the physics edition and
the sum of the M and T readings both give \(w=-\tfrac13\), not the \(w=-1\) of a vacuum term.

\paragraph{Remark v4.10.0-r55.4 (own remainders).}
Independently of the number: existence and self-adjointness of \(D_-\) without the withdrawn Riemann source (OP2); three generations, no ban on a
fourth, Koide not derived (OP3); the \(\mathbb Z_2\) of the neutrinos, \(\delta_{\rm CP}\), CKM and the assignment of gauge charges with multiplicity \(n\) (OP4); the absolute value
of \(\Lambda\) (OP5); the null witness of C5 (form (iii), r53.7); the global closure C6/T2, which the guards do not claim.

\paragraph{Remark v4.10.0-r55.5 (where the open points of the articles sit).}
The three open points of the German article (a native physical operator, a derived assignment of gauge charges, an independently founded
geometric source law) sit in Q1a together with the remainder of OP3, in the remainder of OP4, and in Q1b together with Q2, respectively.  The
progress conditions of the physics edition sort the same way: the mass branch (operator and phase selection) into Q1a and OP3, the generation
branch into the remainders of OP3 and OP4, the geometric branch (coframe, local solution) into Q2.

\paragraph{Remark v4.10.0-r55.6 (consequences for the work plan).}
A negative closure of OP5 (\(\Lambda\) as an input) concerns only the remainder of OP5 and depends neither on OP1 nor on Q1 or Q2.  An attractor
dynamics \(\Omega_M\to\beta^2\) sits in Q2 and needs the OP5 identification in addition.  After r54 the place with the widest reach is Q1a: one orientation
datum would bear on the octant sign and on the counting of OP1 at once.  H118 leaves one variant untested, a temporal order of the cycle, under
which the octant phase would be even under time reversal.

\subsection*{v4.10.0-r55 frames}
\label{subsec:v4100-r55-frames}

No quantity of different frames is compared.  The frames M, T, P and the field rest frame \(R\) appear only as recorded in the cited layers.

\subsection*{v4.10.0-r55 register entries}
\label{subsec:v4100-r55-register}

\begin{center}
\small
\begin{tabular}{@{}>{\raggedright\arraybackslash}p{0.34\textwidth}>{\raggedright\arraybackslash}p{0.58\textwidth}@{}}
\toprule
statement & status after v4.10.0-r55\\
\midrule
one number, five identifications & recorded (r55.1), as in r30.6\\
Q1 (readout) & split into Q1a (orientation, open) and Q1b (frame and axis; r54: not supplied by the readout)\\
Q2 (dynamics) & open; I1 the only candidate touching OP1, OP3 and OP5\\
own remainders of OP2--OP5, C5, C6/T2 & open\\
\bottomrule
\end{tabular}
\end{center}

\subsection*{v4.10.0-r55 notes}
\label{subsec:v4100-r55-notes}

\texttt{notes/\allowbreak research/\allowbreak op\_\allowbreak cross\_\allowbreak map\_\allowbreak 2026\_\allowbreak 11\_\allowbreak 23/\allowbreak README.md} (exploration of 23 November 2026),
\texttt{notes/\allowbreak research/\allowbreak op\_\allowbreak root\_\allowbreak map\_\allowbreak 2026\_\allowbreak 11\_\allowbreak 18/} (root map), layers r53 and r54.

\subsection*{v4.10.0-r55 ledger and guard update}
\label{subsec:v4100-r55-ledger-guard}

\[
R_{\rm chk}^{v4100\text{-}r55}=(1,1,0)
\]
for (status recorded; every statement checked against the cited layers; no open problem derived).

The layer spends no guard:
\[
R_{\rm glob}^{(v4100\text{-}r55,+)}=(0,1,1,1,1),
\qquad
\Gamma_{\rm guard}^{v4100\text{-}r55}=(1,1,1,1,1),
\qquad
\beta_4=1/256\ \text{unspent}.
\]
The layer does not imply exact finite \(C5\) zero closure, exact all-readout terminality, \(C6/T2\), a global photon-clock/RH theorem,
or global all-readout closure.  OP1--OP5, C5 and C6/T2 remain open.
